%%%%%%%%%%%%%%%%%%%%%%%%%%%%%%%%%%%%%%%%%%%%%%%%%%%%%%%%%%%%%%%%%%%%%%%%%%%%%%%%
%%  config/fonts.tex
%%  Fonts and Persian typesetting.  **This file must be the last thing loaded in
%%  the preamble**, because xepersian (and the bidi package behind it) has to be
%%  the last package of the document.
%%
%%  The bundled family is Vazirmatn by Saber Rastikerdar (SIL OFL 1.1), in
%%  fonts/Vazirmatn/.  It covers Persian *and* Latin, so one family carries the
%%  whole deck and mixed Persian/English lines stay optically consistent.
%%
%%  Paths are relative to the directory you compile from, i.e. the folder that
%%  holds main.tex.  To use another family, change the four \...font commands
%%  below and nothing else -- the theme never refers to a concrete font file.
%%%%%%%%%%%%%%%%%%%%%%%%%%%%%%%%%%%%%%%%%%%%%%%%%%%%%%%%%%%%%%%%%%%%%%%%%%%%%%%%

% Proper Computer Modern maths (Vazirmatn has no maths glyphs).  Loaded *before*
% xepersian on purpose: no package may come after it.
\usefonttheme[onlymath]{serif}

%%%%%%%%%%%%%%%%%%%%%%%%%%%%%%%%%%%%%%%%%%%%%%%%%%%%%%%%%%%%%%%%%%%%%%%%%%%%%%%%
%% xepersian -- the last package in the preamble
%%%%%%%%%%%%%%%%%%%%%%%%%%%%%%%%%%%%%%%%%%%%%%%%%%%%%%%%%%%%%%%%%%%%%%%%%%%%%%%%
\usepackage{xepersian}

%%%%%%%%%%%%%%%%%%%%%%%%%%%%%%%%%%%%%%%%%%%%%%%%%%%%%%%%%%%%%%%%%%%%%%%%%%%%%%%%
%% Main Persian font
%%
%%  Mapping=parsidigits turns every ASCII digit of the Persian text into a
%%  Persian one, so slide numbers, chapter numbers and outline numbers come out
%%  as ۱ ۲ ۳ without any extra markup.
%%  Persian has no italics; the italic slot is mapped to the Medium weight, which
%%  gives \emph{...} a gentle emphasis instead of a fake slant.
%%%%%%%%%%%%%%%%%%%%%%%%%%%%%%%%%%%%%%%%%%%%%%%%%%%%%%%%%%%%%%%%%%%%%%%%%%%%%%%%
\settextfont[
  Path            = fonts/Vazirmatn/,
  Extension       = .ttf,
  UprightFont     = *-Regular,
  BoldFont        = *-Bold,
  ItalicFont      = *-Medium,
  BoldItalicFont  = *-ExtraBold,
  Mapping         = parsidigits,
  Scale           = 1.0,
]{Vazirmatn}

%%%%%%%%%%%%%%%%%%%%%%%%%%%%%%%%%%%%%%%%%%%%%%%%%%%%%%%%%%%%%%%%%%%%%%%%%%%%%%%%
%% Latin font
%%
%%  Same family, but *without* parsidigits: English runs keep Latin digits, so
%%  "Precision@3" or "GPT-4" look the way they should.
%%%%%%%%%%%%%%%%%%%%%%%%%%%%%%%%%%%%%%%%%%%%%%%%%%%%%%%%%%%%%%%%%%%%%%%%%%%%%%%%
\setlatintextfont[
  Path            = fonts/Vazirmatn/,
  Extension       = .ttf,
  UprightFont     = *-Regular,
  BoldFont        = *-Bold,
  ItalicFont      = *-Medium,
  BoldItalicFont  = *-ExtraBold,
  Ligatures       = TeX,
  Scale           = 1.0,
]{Vazirmatn}

% The family \en{...} switches to.  Keep it in sync with \setlatintextfont.
\newfontfamily\PTlatinfont[
  Path            = fonts/Vazirmatn/,
  Extension       = .ttf,
  UprightFont     = *-Regular,
  BoldFont        = *-Bold,
  ItalicFont      = *-Medium,
  BoldItalicFont  = *-ExtraBold,
  Ligatures       = TeX,
]{Vazirmatn}

% A true semibold, used for \hl{...} and \term{...} instead of a heavy bold.
%
%  Script=Arabic is not optional here.  \settextfont / \setlatintextfont get the
%  Persian defaults from xepersian, but a hand-rolled \newfontfamily does not:
%  without an explicit Arabic script the shaping features never fire and the
%  letters come out unjoined.
\newfontfamily\PTsemibold[
  Path            = fonts/Vazirmatn/,
  Extension       = .ttf,
  Script          = Arabic,
  UprightFont     = *-SemiBold,
  BoldFont        = *-Bold,
  ItalicFont      = *-SemiBold,
  BoldItalicFont  = *-Bold,
  Mapping         = parsidigits,
]{Vazirmatn}

% A light weight for very large display text, should you want it.
\newfontfamily\PTlight[
  Path            = fonts/Vazirmatn/,
  Extension       = .ttf,
  Script          = Arabic,
  UprightFont     = *-Light,
  BoldFont        = *-Medium,
  ItalicFont      = *-Light,
  BoldItalicFont  = *-Medium,
  Mapping         = parsidigits,
]{Vazirmatn}

% Monospace for the PTcode environment.  Left at the LaTeX default (Latin
% Modern Mono), which every distribution has.  Uncomment to upgrade:
%\setlatinmonofont[Scale=MatchLowercase]{DejaVu Sans Mono}

%%%%%%%%%%%%%%%%%%%%%%%%%%%%%%%%%%%%%%%%%%%%%%%%%%%%%%%%%%%%%%%%%%%%%%%%%%%%%%%%
%% Weight tuning for the theme's inline helpers
%%
%%  Now that a real SemiBold exists, \hl{...} and \term{...} can use it: on a
%%  projector a semibold Persian word reads as emphasis, a bold one reads as
%%  shouting.
%%%%%%%%%%%%%%%%%%%%%%%%%%%%%%%%%%%%%%%%%%%%%%%%%%%%%%%%%%%%%%%%%%%%%%%%%%%%%%%%
\setbeamerfont{PT highlight}{family=\PTsemibold,series=\mdseries}
\setbeamerfont{PT term}     {family=\PTsemibold,series=\mdseries}
